\section{Experiments}
\label{sec:experiments}

\subsection{Experimental Setup}

\textbf{Models and evaluation protocol.}
This work studies visual token compression inside document VLMs, rather than training a new OCR foundation model. We evaluate \ours{} on PaddleOCR-VL-1.6, MinerU2-VLM, and dots.ocr. PaddleOCR-VL-1.6 and MinerU2-VLM originally contain crop-based pipelines. To isolate visual token compression inside the VLM, our compression experiments remove the external crop stage and use a single full-image VLM forward pass. Official pipeline results are kept as reference points for end-to-end performance.

For models with a fixed reduction projector or a fixed image-token placeholder layout, we only adapt the input protocol so that placeholders match the compressed visual-token sequence. The projector parameters and the main VLM backbone are otherwise unchanged.

\textbf{Datasets and metrics.}
Experiments are mainly conducted on OmniDocBench v1.6. Parsing quality is measured by TextEdit, Formula CDM, Table TEDS, Table TEDS\textsubscript{s}, Reading Order Edit, and Overall. Efficiency is measured by input visual tokens, retained visual tokens after compression, TFLOPS breakdown, inference speed, and training speed. We also report category-level compression statistics on OmniDocBench v1.6 and Real5-OmniDocBench to analyze content-adaptive behavior.

\textbf{Compared methods.}
We consider three comparison settings. First, we compare with public end-to-end document parsing systems to position the compressed model on the quality--efficiency frontier. Second, we compare compression rules under the same host model and matched token budgets. We first train a no-merge SFT backbone, and then attach IDPrune-style, ToMe-style, and DivPrune-style as training-free compression operators. In contrast, \ours{} starts from the same initialization, enables BlackHole Merge during SFT, and further trains the lightweight RL Policy.

We use the term ``style'' because these baselines cannot be migrated into a unified VLM-only path as line-by-line equivalent implementations: their official versions depend on different architectures, insertion sites, and token-packing layouts. We therefore preserve their core compression principles instead of reproducing full official systems. ToMe-style merges removed tokens into retained tokens according to feature similarity, DivPrune-style performs diversity-based subset selection, and IDPrune-style selects tokens using visual-token importance with diversity suppression.

\subsection{Quality--Efficiency Comparison}

Figure~\ref{fig:performance} compares \ours{} with representative end-to-end document parsing systems. This comparison positions the compressed model on the quality--efficiency frontier. The following tables further separate official pipelines, VLM-only baselines, and BHM-compressed VLMs.

% Preamble:
% \usepackage{booktabs}
% \usepackage{multirow}
% \usepackage{array}
% \newcolumntype{C}[1]{>{\centering\arraybackslash}m{#1}}
% \newcolumntype{L}[1]{>{\raggedright\arraybackslash}m{#1}}

\begin{table*}[t]
\centering
\fontsize{7.8}{9.0}\selectfont
\setlength{\tabcolsep}{3.0pt}
\renewcommand{\arraystretch}{1.15}
\begin{tabular*}{\textwidth}{@{\extracolsep{\fill}}C{0.18\textwidth}C{0.13\textwidth}L{0.29\textwidth}C{0.07\textwidth}C{0.13\textwidth}C{0.055\textwidth}C{0.055\textwidth}@{}}
\toprule
\multicolumn{2}{c}{\textbf{Model}} &
\multicolumn{1}{c}{\textbf{Compression}} &
\multicolumn{2}{c}{\textbf{Vision Tokens}} &
\multicolumn{2}{c}{\textbf{VLM TFLOPS}} \\
\cmidrule(lr){1-2}
\cmidrule(lr){3-3}
\cmidrule(lr){4-5}
\cmidrule(lr){6-7}
\textbf{Backbone} &
\textbf{Arch} &
\textbf{Method\&Stage} &
\textbf{Input} &
\textbf{Merged} &
\textbf{ViT} &
\textbf{Prefill} \\
\midrule

\multirow[c]{6}{0.18\textwidth}{\centering\shortstack[c]{\textbf{PaddleOCR-VL-1.6}\\\textbf{(0.9B)}}}
& Pipeline$^\dagger$
& Img$\rightarrow$Enc. (Crop) + Enc.$\rightarrow$LLM (4$\times$ Down)
& 4954 & 1202 (-75.7\%) & 6.87 & 0.84 \\
& E2E VLM$^\dagger_{\text{w/o Crop}}$
& Enc.$\rightarrow$LLM (4$\times$ Down)
& 4954 & 1239 (-75.0\%) & 7.13 & 0.87 \\
& E2E VLM
& --
& 4954 & 4954 (0.0\%) & 7.13 & 6.17 \\
& E2E VLM
& Enc.$|_{\text{L}3}$ (BHM)
& 4954 & 1511 (-69.5\%) & 2.45 & 1.46 \\
& E2E VLM
& Enc.$|_{\text{L}6}$ (BHM)
& 4954 & 1421 (-71.3\%) & 2.95 & 1.10 \\
& E2E VLM
& Enc.$|_{\text{L}9}$ (BHM)
& 4954 & 1434 (-71.1\%) & 3.54 & 1.05 \\
\midrule

\multirow[c]{5}{0.18\textwidth}{\centering\shortstack[c]{\textbf{MinerU2-VLM}\\\textbf{(0.9B)}}}
& Pipeline$^\dagger$
& Img$\rightarrow$Enc. (Crop) + Enc.$\rightarrow$LLM (Resize)
& 10604 & 10604 (0.0\%) & 12.57 & 8.25 \\
& E2E VLM
& Enc.$|_{\text{L}20}$ (BHM)
& 10604 & ? & -- & -- \\
& E2E VLM$^\dagger_{\text{w/o Crop}}$
& Img$\rightarrow$Enc. (Resize)
& 729 & 729 (0.0\%) & 0.67 & 0.57 \\
& E2E VLM
& Img$\rightarrow$Enc. (Resize)
& 729 & ? & -- & -- \\
& E2E VLM
& Enc.$|_{\text{L}3}$ (BHM)
& 729 & ? & -- & -- \\
\midrule

\multirow[c]{2}{0.18\textwidth}{\centering\shortstack[c]{\textbf{dots.ocr}\\\textbf{(1.7B)}}}
& E2E VLM$^\dagger$
& Enc.$\rightarrow$LLM (4$\times$ Down)
& 21510 & 5377 (-75.0\%) & 211.36 & 21.67 \\
& E2E VLM
& Enc.$|_{\text{L}32}$ (BHM)
& 21510 & -- & -- & -- \\
\bottomrule
\end{tabular*}
\caption{Efficiency comparison on OmniDocBench v1.6. Method with $^\dagger$ denotes the official implementation.}
\label{tab:full_unified_efficiency}
\end{table*}

\begin{table*}[t]
\centering
\fontsize{7.1}{8.3}\selectfont
\setlength{\tabcolsep}{1.2pt}
\renewcommand{\arraystretch}{1.12}
\begin{tabular*}{\textwidth}{@{\extracolsep{\fill}}
>{\centering\arraybackslash}m{0.11\textwidth}
>{\centering\arraybackslash}m{0.11\textwidth}
>{\centering\arraybackslash}m{0.26\textwidth}
>{\centering\arraybackslash}m{0.05\textwidth}
>{\centering\arraybackslash}m{0.08\textwidth}
>{\centering\arraybackslash}m{0.075\textwidth}
>{\centering\arraybackslash}m{0.075\textwidth}
>{\centering\arraybackslash}m{0.08\textwidth}
>{\centering\arraybackslash}m{0.07\textwidth}
@{}}
\toprule
\multicolumn{2}{c}{\textbf{Model}} &
\multicolumn{1}{c}{\textbf{Compression}} &
\multicolumn{6}{c}{\textbf{Parsing Metrics}} \\
\cmidrule(lr){1-2}
\cmidrule(lr){3-3}
\cmidrule(lr){4-9}
\textbf{Backbone} &
\textbf{Arch} &
\textbf{Method\&Stage} &
\textbf{Text}$^{\scriptsize\textbf{Edit}}\downarrow$ &
\textbf{Formula}$^{\scriptsize\textbf{CDM}}\uparrow$ &
\textbf{Table}$^{\scriptsize\textbf{TEDS}}\uparrow$ &
\textbf{Table}$^{\scriptsize\textbf{TEDS-S}}\uparrow$ &
\textbf{R-order}$^{\scriptsize\textbf{Edit}}\downarrow$ &
\textbf{Overall}$\uparrow$ \\
\midrule

\multirow[c]{6}{0.12\textwidth}{\centering\shortstack[c]{\textbf{PaddleOCR-VL}\\\textbf{(0.9B)}}}
& Pipeline$^\dagger$
& Img$\rightarrow$Enc. (Crop) + Enc.$\rightarrow$LLM (4$\times$ Down)
& 0.033 & 97.49 & 94.76 & 97.11 & 0.127 & 96.33 \\
& E2E VLM$^\dagger_{\text{w/o Crop}}$
& Enc.$\rightarrow$LLM (4$\times$ Down)
& 0.072 & 92.22 & 85.01 & 88.24 & 0.168 & 90.06 \\
& E2E VLM
& --
& 0.045 & 94.00 & 88.32 & 91.31 & 0.141 & 92.61 \\
& E2E VLM
& Enc.$|_{\text{L}3}$ (BHM)
& 0.056 & 95.08 & 87.20 & 90.41 & 0.146 & 92.23 \\
& E2E VLM
& Enc.$|_{\text{L}6}$ (BHM)
& 0.045 & 93.77 & 85.93 & 89.33 & 0.139 & 91.75 \\
& E2E VLM
& Enc.$|_{\text{L}9}$ (BHM)
& 0.049 & 95.50 & 87.62 & 91.04 & 0.138 & 92.76 \\
\midrule

\multirow[c]{5}{0.12\textwidth}{\centering\shortstack[c]{\textbf{MinerU2-VLM}\\\textbf{(0.9B)}}}
& Pipeline$^\dagger$
& Img$\rightarrow$Enc. (Crop) + Enc.$\rightarrow$LLM (Resize)
& 0.102* & 82.42 & 81.38 & 84.88 & 0.198 & 84.54 \\
& E2E VLM
& Enc.$|_{\text{L}20}$ (BHM)
& -- & -- & -- & -- & -- & -- \\
& E2E VLM$^\dagger_{\text{w/o Crop}}$
& Img$\rightarrow$Enc. (Resize)
& 0.130* & 77.40 & 75.53 & 79.05 & 0.224 & 79.96 \\
& E2E VLM
& Img$\rightarrow$Enc. (Resize)
& -- & -- & -- & -- & -- & -- \\
& E2E VLM
& Enc.$|_{\text{L}3}$ (BHM)
& -- & -- & -- & -- & -- & -- \\
\midrule

\multirow[c]{2}{0.12\textwidth}{\centering\shortstack[c]{\textbf{dots.ocr}\\\textbf{(1.7B)}}}
& E2E VLM$^\dagger$
& Enc.$\rightarrow$LLM (4$\times$ Down)
& 0.054* & 90.88 & 86.51 & 89.98 & 0.139 & 90.68 \\
& E2E VLM
& Enc.$|_{\text{L}32}$ (BHM)
& -- & -- & -- & -- & -- & -- \\
\bottomrule
\end{tabular*}
\caption{Parsing performance comparison on OmniDocBench v1.6. Method with $^\dagger$ denotes the official implementation.}
\label{tab:full_unified_parsing}
\end{table*}


Table~\ref{tab:full_unified_efficiency} and Table~\ref{tab:full_unified_parsing} show that fixed reduction in official pipelines mainly reduces tokens after the visual encoder, whereas BHM reduces tokens inside the encoder. For PaddleOCR-VL-1.6, BHM reduces the visual sequence from 4954 tokens to about 1400 retained tokens, lowering both the later ViT cost and LLM prefill cost while maintaining competitive parsing quality.

\subsection{Controlled Compression Comparison}
\label{sec:controlled_comparison}

\begin{table}[t]
\centering
\fontsize{7.6}{8.8}\selectfont
\setlength{\tabcolsep}{2.6pt}
\renewcommand{\arraystretch}{1.15}
\begin{tabular}{@{}lccccc c@{}}
\toprule
\textbf{Method} &
\shortstack{\textbf{Text}\\\textbf{Edit}$\downarrow$} &
\shortstack{\textbf{Formula}\\\textbf{CDM}$\uparrow$} &
\shortstack{\textbf{Table}\\\textbf{TEDS}$\uparrow$} &
\shortstack{\textbf{Table}\\\textbf{TEDS-S}$\uparrow$} &
\shortstack{\textbf{R-order}\\\textbf{Edit}$\downarrow$} &
\textbf{Overall}$\uparrow$ \\
\midrule
\multicolumn{7}{@{}l}{\textit{Baseline Methods}} \\
DivPrune-Style     & 0.990 & 0.48 & 0.00 & 0.00 & 0.768 & 0.32 \\
IDPrune-Style      & 0.996 & 0.00 & 0.00 & 0.00 & 0.805 & 0.13 \\
ToMe-Style         & 0.683 & 0.00 & 0.81 & 12.83 & 0.566 & 15.50 \\
\midrule
\multicolumn{7}{@{}l}{\textit{Ours: BlackHoleMerge}} \\
BlackHoleMerge-L3  & 0.056 & 95.08 & 87.20 & 90.41 & 0.146 & 92.23 \\
BlackHoleMerge-L6  & 0.045 & 93.77 & 85.93 & 89.33 & 0.139 & 91.75 \\
BlackHoleMerge-L9  & 0.049 & 95.50 & 87.62 & 91.04 & 0.138 & 92.76 \\
\bottomrule
\end{tabular}
\caption{Controlled comparison between training-free compression styles and BlackHole Merge on OmniDocBench v1.6.}
\label{tab:omnidocbench_compare}
\end{table}

Table~\ref{tab:omnidocbench_compare} isolates the compression rule under a controlled VLM-only setting. The style baselines are attached to the same no-merge SFT backbone and matched to a similar token budget, while \ours{} uses merge-aware SFT and RL from the same initialization. Under this protocol, BlackHole Merge preserves OCR quality much better than training-free pruning or merging styles, showing that adapting the backbone to the compression path is important for document parsing.

\subsection{Compression Behavior Analysis}

\begin{table}[t]
\centering
\fontsize{7.4}{8.6}\selectfont
\setlength{\tabcolsep}{2.0pt}
\renewcommand{\arraystretch}{1.12}
\begin{tabular*}{\columnwidth}{@{\extracolsep{\fill}}%
>{\centering\arraybackslash}m{0.13\columnwidth}
>{\centering\arraybackslash}m{0.38\columnwidth}
>{\centering\arraybackslash}m{0.13\columnwidth}
>{\centering\arraybackslash}m{0.13\columnwidth}
>{\centering\arraybackslash}m{0.13\columnwidth}
@{}}
\toprule
\multicolumn{2}{c}{\textbf{Experiment}} &
\multicolumn{3}{c}{\textbf{Token Density}} \\
\cmidrule(lr){1-2}
\cmidrule(lr){3-5}
\shortstack[c]{\textbf{Merge}\\\textbf{Layer}} &
\textbf{Strategy} &
\textbf{Text*} &
\textbf{Blank} &
\textbf{Gap*} \\
\midrule
-- & Pretrained & 0.597 & 0.441 & 0.156 \\
\midrule
\multirow[c]{2}{*}{\textbf{Layer 3}}
& Merge-aware SFT & 0.642 & 0.388 & 0.254 \\
& Merge SFT + Confidence-Penalty RL & \textbf{0.644} & \textbf{0.071} & \textbf{0.572} \\
\cmidrule(lr){2-5}
\multirow[c]{2}{*}{\textbf{Layer 6}}
& Merge-aware SFT & \textbf{0.625} & 0.355 & 0.270 \\
& Merge SFT + Confidence-Penalty RL & 0.562 & \textbf{0.051} & \textbf{0.512} \\
\cmidrule(lr){2-5}
\multirow[c]{2}{*}{\textbf{Layer 9}}
& Merge-aware SFT & \textbf{0.617} & 0.423 & 0.194 \\
& Merge SFT + Confidence-Penalty RL & 0.565 & \textbf{0.082} & \textbf{0.483} \\
\bottomrule
\end{tabular*}
\caption{Token density evolution across different merge layers and training stages. Text/Blank metrics are computed on text and blank/background regions derived from OmniDocBench bounding boxes.}
\label{tab:token_density}
\end{table}

\begin{table}[t]
\centering
\fontsize{7.4}{8.6}\selectfont
\setlength{\tabcolsep}{2.0pt}
\renewcommand{\arraystretch}{1.12}
\begin{tabular*}{\columnwidth}{@{\extracolsep{\fill}}%
>{\centering\arraybackslash}m{0.13\columnwidth}
>{\centering\arraybackslash}m{0.36\columnwidth}
>{\centering\arraybackslash}m{0.11\columnwidth}
>{\centering\arraybackslash}m{0.10\columnwidth}
>{\centering\arraybackslash}m{0.10\columnwidth}
>{\centering\arraybackslash}m{0.10\columnwidth}
@{}}
\toprule
\multicolumn{2}{c}{\textbf{Experiment}} &
\multicolumn{4}{c}{\textbf{Keep Ratio}} \\
\cmidrule(lr){1-2}
\cmidrule(lr){3-6}
\shortstack[c]{\textbf{Merge}\\\textbf{Layer}} &
\textbf{Strategy} &
\shortstack[c]{\textbf{Gray}\\\textbf{Zone}} &
\textbf{Text*} &
\textbf{Blank} &
\textbf{Overall} \\
\midrule
-- & Pretrained & 68.7\% & 77.1\% & 41.7\% & 57.0\% \\
\midrule
\multirow[c]{2}{*}{\textbf{Layer 3}}
& Merge-aware SFT & 51.7\% & 78.7\% & 23.8\% & 47.4\% \\
& Merge SFT + Confidence-Penalty RL & \textbf{2.4\%} & 64.7\% & \textbf{6.9\%} & \textbf{31.8\%} \\
\cmidrule(lr){2-6}
\multirow[c]{2}{*}{\textbf{Layer 6}}
& Merge-aware SFT & 48.0\% & 79.0\% & 21.5\% & 45.8\% \\
& Merge SFT + Confidence-Penalty RL & \textbf{3.5\%} & 56.3\% & \textbf{4.1\%} & \textbf{26.6\%} \\
\cmidrule(lr){2-6}
\multirow[c]{2}{*}{\textbf{Layer 9}}
& Merge-aware SFT & 61.9\% & 80.6\% & 37.0\% & 55.8\% \\
& Merge SFT + Confidence-Penalty RL & \textbf{4.2\%} & 56.5\% & \textbf{6.6\%} & \textbf{28.1\%} \\
\bottomrule
\end{tabular*}
\caption{Keep ratio statistics across different merge layers and training stages. Gray Zone indicates the ambiguous region between text and blank/background tokens.}
\label{tab:keep_ratio}
\end{table}

Table~\ref{tab:merged_density_and_penalty} analyzes how merge-aware SFT and RL shape token densities. Merge-aware SFT increases the density gap between text and background tokens, indicating that the model adapts its Key-norm ranking to the compressed visual path. Confidence-Penalty RL further polarizes the distribution, suppresses blank-token density, and reduces the gray zone, turning the continuous SFT ranking into a clearer keep/drop decision.

\begin{table}[!h]
  \fontsize{8.5}{10.0}\selectfont
  \setlength{\tabcolsep}{6pt}
  \renewcommand{\arraystretch}{1.2}
\centering
\begin{tabular}{>{\raggedright\arraybackslash}m{0.54\linewidth}|>{\centering\arraybackslash}m{0.39\linewidth}}
\toprule
\multicolumn{1}{c|}{\textbf{Benchmark / Sub-Category}} & \textbf{Avg. Token Compression Rate ($\uparrow$)} \\
\midrule
\multicolumn{2}{l}{\textit{OmniDocBench v1.6 Split}} \\
PPT       & 79.82\% (\textit{+17.97\%}) \\
Book      & 75.20\% (\textit{+10.83\%}) \\
Paper     & 64.49\% (\textit{+10.27\%}) \\
Exam      & 65.32\% (\textit{+4.52\%})  \\
Textbook  & 75.57\% (\textit{+5.53\%})  \\
Newspaper & 41.17\% (\textit{+0.52\%})  \\
Magazine  & 56.81\% (\textit{-16.32\%}) \\
Report    & 74.89\% (\textit{+18.09\%}) \\
Note      & 74.12\% (\textit{+22.79\%}) \\
\midrule
\multicolumn{2}{l}{\textit{Real5-OmniDocBench Split}} \\
Scanning               & 71.36\% (\textit{+26.93\%}) \\
Warping  & 78.36\% (\textit{+33.93\%}) \\
Screen-Photography & 78.07\% (\textit{+33.64\%}) \\
Illumination & 75.55\% (\textit{+31.12\%}) \\
Skew & 76.14\% (\textit{+31.71\%}) \\
\bottomrule
\end{tabular}
   \caption{Average visual-token compression rates across document sub-categories on OmniDocBench v1.6 and Real5-OmniDocBench. Values in parentheses denote improvements over the corresponding reference compression bound.}
   \label{tab:omnidocbench_compression}
\end{table}

Table~\ref{tab:omnidocbench_compression} further reports category-level compression behavior. The learned density is content-adaptive rather than tied to a fixed page layout: sparse or background-heavy categories generally obtain larger compression gains, while text-heavy layouts retain more tokens. Qualitative visualizations show the same trend, with text, formulas, tables, and layout boundaries preferentially retained and blank regions aggressively merged.